\section{From fixed blocks to cutoff}
\label{sec:cutoff}

We now convert the fixed-time spectral estimates into cutoff.  Uniform
contraction on one sufficiently long fixed block gives the upper bound.
For the lower bound, we use a linear statistic if the rate is the spin
edge, and prepare a slow visible mode otherwise.  Throughout,
$0<\beta<1$ is fixed, $\lambda$ is the deterministic number in
\Cref{eq:exact-lambda}, and Euclidean norms remain unnormalized.  For a forward trajectory $(X_s)$ and a terminal
function $f$ at time $T$, write
\[
 p_s(x)=dS_{T-s}f(x),\qquad
 N_s^2=\E\norm{p_s(X_s)}^2,\qquad 0\le s\le T.
\]
Thus $p$ evolves backwards while the law in the expectation evolves
forwards.  Bounds stated uniformly for polynomial tests mean that, for
each fixed $M<\infty$, they hold for every real $f$ with
$\norm f_\infty\le n^M$.  Increasing $M$ by a fixed amount is always
allowed before selecting the good-disorder event.  In particular,
normalization by a norm at least an inverse power of $n$ preserves this
class.

\subsection{Uniform blocks and deterministic selection}

\paragraph{Restarting on a fixed interval.}
The limiting statements of \Cref{sec:exact} use prescribed comparison
times.  We will apply them to blocks whose positions may depend on the
dimension and the disorder.  The following finite-chain identities put
such a block on a fixed relative interval before any limit is taken.
Fix its length $\ell>0$.  Conditional on $J$, let the original chain
start from $\nu$, and choose $0\le a\le T-\ell$ and a terminal
function $f$.  These choices may depend on $n,J$, but not on the
realized trajectory.  Set
\[
 Y_r=X_{a+r},\qquad \widehat\nu=\nu S_a,\qquad
 g=S_{T-a-\ell}f,\qquad 0\le r\le\ell.
\]
For $0=r_0<r_1<\cdots<r_k\le\ell$, the joint law is
\[
 \P(Y_{r_0}=x_0,\ldots,Y_{r_k}=x_k)
 =\widehat\nu(x_0)\prod_{j=1}^k
      S_{r_j-r_{j-1}}(x_{j-1},x_j).
\]
Thus $Y$ is the same ordinary heat-bath chain with initial law
$\widehat\nu$.  Its actual forward and reverse conditional laws,
and trees copied at a single endpoint, are those of the original
path restricted to the block.  The semigroup identity gives
\begin{equation}\label{eq:cutoff-restart}
 \widehat p_r:=dS_{\ell-r}g=p_{a+r},\qquad
 \E_{\widehat\nu S_r}\norm{\widehat p_r}^2=N_{a+r}^2.
\end{equation}
It preserves exactness, the Hessian budgets, and all contractions with
fixed recipe vectors at corresponding endpoints.

The two bounds needed for the limiting arguments are also preserved.
Positivity of $S_t$ gives $\norm g_\infty\le\norm f_\infty$ and
$\sup_x\norm{\widehat p_r(x)}\le\sqrt n\,\norm f_\infty$;
stationarity also preserves the equilibrium mean, $\mu g=\mu f$.
Division by an endpoint norm bounded below by an inverse power of $n$
therefore changes only the fixed polynomial envelope.  If
$h=\dd\nu/\dd\mu$, reversibility and positivity give
\begin{equation}\label{eq:cutoff-restart-density}
 \frac{\dd(\nu S_a)}{\dd\mu}=S_a h,
 \qquad \norm{S_a h}_\infty\le\norm h_\infty.
\end{equation}
At relative time $r$, its marginal density is
$\widehat h_r=S_r\widehat h_0=h_{a+r}$.  For $0\le r<u\le\ell$,
the reverse conditional law is therefore
\[
 \P(Y_r=x\mid Y_u=y)
 =\frac{\mu(x)\widehat h_r(x)S_{u-r}(x,y)}
        {\mu(y)\widehat h_u(y)}.
\]
In particular, the edge density ratios in the reverse Hellinger
budget are unchanged at corresponding times.  An upper relative
log-density $o(n)$ survives every restart.  For a point start and $a\ge b$, the same argument applied
after time $b$ gives
$\norm{h_a}_\infty\le\norm{h_b}_\infty$; hence the fixed-$b$
bound in \Cref{prop:stat-density} controls all later block origins
on its single disorder event.  A trajectory Poincar\'e bound, when
assumed, has exactly the same constant and remainder after this
change of coordinates because the marginal laws are identical.

We will use $\ell=D$ or $2D$ for a fixed block length $D$.
The endpoints $0,D$ (and $2D$ when $\ell=2D$) can then be retained as prescribed times in
\Cref{prop:exact-transport}.  Interior comparison times are still
selected using \Cref{lem:exact-good-times} and the integrated
comparison bounds.  The restart does not replace that selection by a
supremum over exceptional auxiliary times.

\begin{proposition}[Uniform contraction on a fixed block]
\label{prop:cutoff-block}
Fix $0<\xi<\lambda$, a polynomial envelope $M$, a horizon constant
$A<\infty$, and $H<\infty$.  There are deterministic constants
$D,b_0<\infty$ and disorder events of probability tending to one on
which the following assertion holds simultaneously for all starting
configurations, all such terminal functions, and all
$b_0\le a<a+D\le T\le A\log n$:
\begin{equation}
 N_a^2\le e^{-2(\lambda-\xi)D}N_{a+D}^2+n^{-H}.
 \label{eq:cutoff-block}
\end{equation}
The same assertion, with no burn-in, holds in the deterministic
typicality convention of \Cref{thm:path-calculus} for families of
initial laws satisfying
$\log\norm{\dd\nu/\dd\mu}_\infty=o(n)$ uniformly.
The constants may be chosen after $H$ and $M$; no quantitative rate
for a fixed-time $o(1)$ is asserted or required.
\end{proposition}

\begin{proof}
The proof has three steps: a limiting contraction with strict slack,
uniformity down to an inverse-polynomial norm, and a fixed burn-in.
Begin with a block whose initial law has vanishing specific log density.
Work first with a fixed $D>2$, which will be enlarged below.
For each possible choice of its origin and terminal horizon, apply
\Cref{eq:cutoff-restart} with $\ell=D$ and terminal function
$g=S_{T-a-D}f$.  In the limiting argument we use these restarted
objects, denoting the process again by $X$, so that $a=0$, $T=D$,
and the endpoint times are prescribed.  Normalize $N_D=1$.
The rough square estimate in \Cref{prop:path-rough}, integrated over
the last unit of this interval, supplies a time
$s\in[D-1,D-1/2]$ at which
\begin{equation}
 \E\norm{p_s}^2+
 \E\sum_jw_j\norm{d_jp_s}^2\le C.
 \label{eq:cutoff-good-time}
\end{equation}
One can impose the reverse Hellinger estimate of
\Cref{lem:exact-good-times} at the same time.  Both requirements hold
on a set of times of positive measure after enlarging their fixed
constants.  By \Cref{prop:exact-closedness}, the visible projection
$q_s$ therefore lies in the exact weighted sector, including its
scalar Hilbert factor.  By
\Cref{prop:exact-transport,thm:exact-sector,lem:exact-upgrade}, its
propagation to the prescribed time $0$ has norm
at most
\[
 C_\rho e^{-\rho(D-1)},\qquad \rho<\lambda.
\]
Conditional expectation on the scalar factor is a contraction.  The
upgrade applies to an unweighted visible vector whose weighted image
lies in $\mathcal X_v$; it does not require membership in
$\mathcal X$.  The hidden remainder at time $0$ satisfies the corresponding estimate
with any rate strictly below $\lambda_{\rm sp}$, by
\Cref{thm:hidden-propagation}, testing against that remainder itself.
Finite recipe projections followed by increasing the projection rank
justify this test; the norm and Hessian bounds in
\Cref{eq:cutoff-good-time} are exactly its input hypotheses.
Consequently, for every $\rho<\lambda$,
\begin{equation}
 \limsup N_0\le C_\rho e^{-\rho(D-1)}.
 \label{eq:cutoff-limiting-block}
\end{equation}
Choose $\rho\in(\lambda-\xi,\lambda)$ and then $D$ so large that the
right side is strictly smaller than $e^{-(\lambda-\xi)D}$.
This leaves fixed multiplicative slack.

We specify how the comparison time is selected, since the path estimates
also contain auxiliary integration variables.  First impose finitely many
recipe, conditional-pair, and matrix tests at a fixed accuracy.  The exponential stationary
discard bounds transfer through the upper density bound and every
fixed polynomial weight.  Bounds for deterministic masks are
integrated over their simplex variables before this transfer.  Fubini
and extraction give convergence in measure on the bounded domains of
auxiliary times.  The uniform rough $L^2$ estimates, with the Hessian
budget when it occurs, control the resulting integrals.  In
particular, we do not bound an exceptional time integral by its
Lebesgue measure times a polynomial pointwise supremum.  Increasing
the finite test lists and decreasing their accuracies permits one
common choice of $s$ along a subsequence.  Strong continuity, or a
further fixed positive smoothing time, treats a convergent sequence
of selected times.  This is the deterministic countable-test content
of \Cref{thm:path-calculus}.

To obtain the additive accuracy in \Cref{eq:cutoff-block}, consider an
arbitrary proposed violating sequence of starts, functions, origins
$a_n$, and horizons $T_n$.  Restart each at $a_n$ on $[0,D]$, with
$g_n=S_{T_n-a_n-D}f_n$, as above, and set $B=N_D^2$ for its
restarted endpoint norm.  If
\[
 B\le e^{-CD}n^{-H},
\]
the rough estimate already gives $N_0^2\le n^{-H}$.
Otherwise divide $p$ and its potential by $\sqrt B$.  A polynomial
pointwise bound remains, since $B^{-1/2}\le e^{CD/2}n^{H/2}$.
If the claimed inequality failed along a sequence of deterministic
good disorders, its normalized version would contradict the strict
slack in \Cref{eq:cutoff-limiting-block}.  Indeed the restarted laws
satisfy \Cref{eq:cutoff-restart-density}, their normalized terminal
functions remain polynomially bounded, and every endpoint comparison
is now at $0$ or $D$.  Only the interior good time is selected after
imposing the integrated tests.  Equation~\eqref{eq:cutoff-restart}
then returns the estimate to the original block, regardless of the
sequences $a_n,T_n$.  This proves the required uniformity.  The
inverse-power remainder comes from normalization and strict slack,
rather than a rate assigned to a fixed-time $o(1)$.

We finally prove that one \emph{fixed} burn-in works with disorder
probability tending to one.  A deterministic diagonal contradiction
will place every required typicality test on the same sequence of
disorders; convergence in probability of the limiting block estimate
alone would not supply this conclusion.  Enumerate the
deterministic good-disorder inequalities just used, including the
quenched exponential bounds for stationary exceptional sets and the
uniform forward superpolynomial discards.  For the first $k$ tests,
with their fixed tolerances, let their positive exponential rates be
$c_1,\ldots,c_k$.  By \Cref{prop:stat-density}, choose fixed numbers
$b_k\uparrow\infty$ so that
\[
 \delta_k=C e^{-c b_k^{2/3}}\longrightarrow0,
 \qquad
 \delta_k<\tfrac14\min_{j\le k}c_j.
\]
Include the density bound at this particular $b_k$ in the finite
good-disorder event $G_{n,k}$.

Suppose every fixed burn-in failed.  For each $b_k$, there would be
$\varepsilon_k>0$ and arbitrarily large dimensions at which its
failure event $F_{n,b_k}$ had probability at least
$\varepsilon_k$.  Choose one such $n_k$, increasing sufficiently
fast, with
\[
 \P(G_{n_k,k}^{\mathrm c})<\varepsilon_k/2,
\]
and with all the finitely many dimension thresholds satisfied.  Select
a deterministic disorder in $F_{n_k,b_k}\cap G_{n_k,k}$, together
with a violating start, function $f_k$, origin $a_k\ge b_k$, and
horizon $T_k\ge a_k+D$.  Restart on $[0,D]$ with initial law
$S_{a_k}(x_k,\cdot)$ and terminal function
$g_k=S_{T_k-a_k-D}f_k$.  Its initial density is at most
$e^{n_k\delta_k}$ by \Cref{eq:cutoff-restart-density} applied after
$b_k$, and $\norm{g_k}_\infty\le\norm{f_k}_\infty$.
Normalize only above the inverse-polynomial cutoff specified above.  Every fixed
typicality test now holds eventually on this deterministic sequence;
its exponential discard survives multiplication by
$e^{n_k\delta_k}$ and the fixed polynomial normalization.  The
limiting argument therefore contradicts the same fixed contraction
slack.  No lower bound on the sequence $\varepsilon_k$ was used.
Thus some fixed $b_0$ works, as asserted.
\end{proof}

\begin{proposition}[Propagation over logarithmic times]
\label{prop:cutoff-gradient}
For $0<\xi<\lambda$ and fixed $A,M,H$, on disorder events of
probability tending to one,
\begin{equation}
 \E\norm{p_a(X_a)}^2
 \le C_\xi e^{-2(\lambda-\xi)(b-a)}
             \E\norm{p_b(X_b)}^2+C_\xi n^{-H},
 \qquad 0\le a\le b\le T\le A\log n.
 \label{eq:cutoff-propagation}
\end{equation}
The estimate is uniform over arbitrary initial laws and all polynomial
tests in the specified envelope.  For every deterministic $u\in\R^n$
with $\norm u=1$, it implies the pointwise envelope
\begin{equation}
 \sup_x\norm{dS_t(u^tx)(x)}^2
 \le C_\xi e^{-2(\lambda-\xi)t}+C_\xi n^{-H},
 \qquad 0\le t\le A\log n.
 \label{eq:cutoff-linear-envelope}
\end{equation}
\end{proposition}

\begin{proof}
Iterate the uniform block estimate after the fixed burn-in.  Contraction
makes the additive remainders summable: for $k$ full blocks their total
is at most
\[
 n^{-H}\sum_{j=0}^{k-1}e^{-2(\lambda-\xi)Dj}
 \le\frac{n^{-H}}{1-e^{-2(\lambda-\xi)D}}.
\]
The rough square estimate handles the burn-in and the remaining
interval of length less than $D$.  Their total length is bounded by
$b_0+D$, so they change only the constant.  One may restart this
argument at time $a$, with its arbitrary marginal law; integrating
the point-start estimate handles mixtures.  Short intervals are
covered directly by the same rough estimate.  Increase the chosen
block accuracy by one fixed power if necessary to absorb constant
factors.  This proves \Cref{eq:cutoff-propagation}.  For $f(x)=u^tx$,
the terminal gradient is the constant vector $u$, giving
\Cref{eq:cutoff-linear-envelope}.
\end{proof}

Two consequences will be used repeatedly: bounded exponential moments
for linear tests, and a common unweighted Poincar\'e bound along evolved
laws.  Write $\rho=\lambda-\xi>0$.  For a unit linear test, the function
\[
 b(r)=C_\xi e^{-2\rho r}+C_\xi n^{-H}
\]
is a deterministic pointwise bound for its squared unweighted
gradient.  Both $\sup b$ and $\int_0^{A\log n}b(r)\dd r$ are bounded
independently of $n$, provided $H>0$.  \Cref{lem:input-mgf} consequently gives
\begin{equation}
 \sup_{\norm u=1}\sup_{x,\,t\le A\log n}
 \E_x e^{\theta(u^tX_t-S_t(u^tx)(x))}\le C_\theta,
 \qquad \theta\in\R.
 \label{eq:cutoff-linear-mgf}
\end{equation}
Indeed its exponent is
$2\theta^2e^{2|\theta|\sqrt{\sup b}}\int b$.
The application uses the unweighted \emph{pointwise} envelope in
\Cref{eq:cutoff-linear-envelope}, as required by the finite-chain moment
lemma.  In particular, all these transition variances are bounded.

The second consequence concerns an unweighted Poincar\'e inequality
along a trajectory.  Suppose
\[
 \Var_\nu g\le C_0\,\nu\norm{dg}^2.
\]
Total variance, the martingale variance identity, and
\Cref{eq:cutoff-propagation} give, for polynomial $f$ and
$t\le A\log n$,
\begin{equation}
 \Var_{\nu S_t}f
 \le C_\rho(C_0+2/\rho)\,\nu S_t\norm{df}^2+O(n^{-H}).
 \label{eq:cutoff-trajectory-poincare}
\end{equation}
To check the constants, the martingale contribution is at most
$4\int_0^t\E\norm{dS_{t-s}f(X_s)}^2\dd s$ and the initial
contribution is at most $C_0\nu\norm{dS_tf}^2$.
The exponential integral is at most $1/(2\rho)$.
Apply \Cref{eq:cutoff-propagation} with an accuracy two powers better
than the displayed one to absorb the time factor.  Thus the constant
in \Cref{eq:cutoff-trajectory-poincare} is fixed independently of the
test's particular polynomial normalization.

\subsection{The upper bound and the spin-edge lower bound}

\begin{lemma}[Euclidean oscillation bound]
\label{lem:cutoff-oscillation}
For every real function $g$ on the cube,
\begin{equation}
 \max g-\min g\le2\sqrt n\sup_x\norm{dg(x)}.
 \label{eq:cutoff-oscillation}
\end{equation}
\end{lemma}

\begin{proof}
Extend $g$ multilinearly to the cube: let $\overline g(z)$,
$z\in[-1,1]^n$, be its expectation under the product law with coordinate
means $z_i$.  Then
$\partial_i\overline g(z)=\E_z d_ig$, since $d_ig$ does not depend
on its own spin.  Jensen's inequality gives
\[
 \norm{\nabla\overline g(z)}^2
 \le\E_z\sum_i(d_ig)^2
 \le\sup_x\norm{dg(x)}^2.
\]
Integrate along the Euclidean segment joining two vertices, whose
length is at most $2\sqrt n$.
\end{proof}

By spin symmetry the equilibrium mean of a linear test is zero.
\Cref{eq:cutoff-linear-envelope,eq:cutoff-oscillation} therefore imply
\begin{equation}
 |S_t(u^tx)(x)|
 \le C\sqrt n\bigl(e^{-\rho t}+n^{-H/2}\bigr),
 \qquad \norm u=1.
 \label{eq:cutoff-linear-mean}
\end{equation}
Put $t_n=\log n/(2\lambda)$.  Fix $\epsilon>0$, and choose $\xi>0$
so small that $\rho(1+\epsilon)>\lambda$.  Since the row norms of
$J$ are bounded, \Cref{eq:cutoff-linear-mgf,eq:cutoff-linear-mean}
give, for every $\delta>0$ and every fixed $R<\infty$,
\begin{equation}
 \sup_x\sup_{(1+\epsilon)t_n\le s\le A\log n}
 \P_x\!\left(\max_i|(JX_s)_i|>\delta\log n\right)
 \le n^{-R}
 \label{eq:cutoff-field-tail}
\end{equation}
for all sufficiently large $n$.  For clarity, choose a fixed
exponential parameter larger than $(R+2)/\delta$, with the bounded
row-norm factor included, use that every field mean is $o(1)$, and
then sum the $n$ bounds.  The estimate is uniform in $s$; no assertion
about a maximum over a continuum of trajectory times is needed.

\begin{lemma}[Upper cutoff]
\label{lem:cutoff-upper}
For every $\epsilon>0$,
\[
 \sup_x\norm{S_{(1+\epsilon)t_n}(x,\cdot)-\mu}_{\TV}
 \longrightarrow0
\]
in disorder probability.
\end{lemma}

\begin{proof}
We will find a late time with bounded gradient energy, propagate that
bound back to time zero, and then control oscillation.  Relabelling
$\epsilon$, it suffices to work at $T=(1+2\epsilon)t_n$.  Fix
$|f|\le1$ and a starting point.  The backward martingale
$S_{T-s}f(X_s)$ has jump bracket density
$2\sum_iw_i p_{s,i}^2$.  Hence
\begin{equation}
 \int_0^T\E\sum_iw_i p_{s,i}^2\dd s\le\tfrac12.
 \label{eq:cutoff-energy-budget}
\end{equation}
On the interval $I=[(1+\epsilon)t_n,T]$, whose length is
$\epsilon t_n$, \Cref{eq:cutoff-field-tail} applies.  On its good
field event,
\[
 w_i\ge\frac{2}{1+e^{2|(JX_s)_i|}}\ge n^{-2\delta}.
\]
On the complementary event $\norm{p_s}^2\le n$, since $|f|\le1$.
There is thus a time $s\in I$ at which
\begin{equation}
 N_s^2\le\frac{n^{2\delta}}{2\epsilon t_n}+n^{1-R}.
 \label{eq:cutoff-terminal-energy}
\end{equation}
Apply \Cref{eq:cutoff-propagation} between $0$ and this time.  Uniformly
over the starting point and $f$,
\[
 \norm{dS_Tf(x)}^2
 \le C n^{-(1+\epsilon)\rho/\lambda+2\delta}
       +C n^{1-R}+C n^{-H}.
\]
Choose
\[
 0<2\delta<(1+\epsilon)\rho/\lambda-1,
 \qquad R>3,\quad H>2.
\]
This choice makes the squared gradient $o(n^{-1})$, exactly the scale
needed to offset the cube diameter in \Cref{lem:cutoff-oscillation}.
Hence $\operatorname{osc}(S_Tf)=o(1)$ uniformly for $|f|\le1$.  Since $\mu S_Tf=\mu f$, the variational
formula for total variation proves the result.
\end{proof}

\begin{lemma}[Lower cutoff when the spin edge determines the rate]
\label{lem:cutoff-spin-lower}
If $\lambda=\lambda_{\rm sp}$, then for every $0<\epsilon<1$,
\[
 \sup_x\norm{S_{(1-\epsilon)t_n}(x,\cdot)-\mu}_{\TV}
 \longrightarrow1
\]
in disorder probability.
\end{lemma}

\begin{proof}
Let $\sigma_n$ be the probability measure characterized by
\[
 \frac1n\sum_i\ip{x_i}{S_tx_i}_\mu
       =\int e^{-rt}\sigma_n(\dd r).
\]
Its fixed-test weak convergence to the spin spectral law is supplied
by \Cref{prop:spin-spectrum}.  For each fixed $\delta>0$, the
definition of its support edge supplies $c_\delta>0$ such that,
with disorder probability tending to one,
\[
 \sigma_n([0,\lambda+\delta])\ge c_\delta.
\]
One obtains this statement by testing a nonnegative continuous
function supported below $\lambda+\delta$ with positive limiting
integral.  It follows for \emph{every} $t\ge0$ that
\begin{equation}
 \E_\mu[X_0^tX_t]\ge c_\delta n e^{-(\lambda+\delta)t}.
 \label{eq:cutoff-spin-overlap}
\end{equation}
The fixed spectral-mass bound remains valid for every time in this
inequality.  It therefore permits logarithmic times without taking a
dimension-dependent Laplace-transform limit.

Averaging over the initial equilibrium spin selects a configuration $x$
for which $\E_x[x^tX_t]$ is at least the value in
\Cref{eq:cutoff-spin-overlap}.  The unit linear test
$f_x(y)=n^{-1/2}x^ty$ detects this remaining overlap.  Its equilibrium mean is zero and its
transition mean is at least
$c_\delta\sqrt n e^{-(\lambda+\delta)t}$.
Its transition variance is bounded by
\Cref{eq:cutoff-linear-mgf}; its equilibrium variance is at most the
inverse gap, by \Cref{thm:input-gap}.  At
$t=(1-\epsilon)t_n$, choose, for example,
$0<\delta<\lambda\epsilon/4$.  The displayed mean diverges as a
positive power of $n$.  Chebyshev's inequality applied at half this
mean gives total variation tending to one.
\end{proof}

\subsection{Alignment with an isolated visible eigenspace}

Assume henceforth that $\lambda<\lambda_{\rm sp}$.
\Cref{prop:exact-isolation} supplies a finite-dimensional slow
eigenspace.  Write $E=(E_1,\ldots,E_m)$ for a basis orthonormal in
the weighted metric.  Its columns have finite unweighted norm; let
$C_E$ bound the norm of the embedding $c\mapsto Ec$ into that space.
The projection onto this eigenspace is the same under the weighted
spectral decomposition and the unweighted isolated spectral
decomposition.  All its complements needed below decay at a rate
strictly greater than $\lambda$.

We approximate the slow columns twice, for different purposes.  Coarse
smooth recipes $B$ control coefficient fluctuations; their derivative
constants are fixed before choosing the alignment block.  Fine recipes
$\widetilde E$ then report the coefficients and define the preparation.
Keeping these choices separate prevents the finer derivative constants
from changing the block length that has already been chosen.

\begin{proposition}[Mean and fluctuation alignment]
\label{prop:cutoff-coefficients}
Fix a logarithmic horizon, a polynomial envelope, and a common
constant $C_P$ in the defective trajectory Poincar\'e bound
\Cref{eq:cutoff-trajectory-poincare}, with arbitrarily prescribed
inverse-power accuracy.  Consider uniformly typically initialized
families of ordinary trajectories satisfying this bound.
For every $\eta>0$, one can choose a fixed block length $D$ and fixed
smooth recipes $\widetilde E$ such that, on disorder events of
probability tending to one, the following holds uniformly.
Put
\[
 A_r=\E\bigl[\widetilde E(X_r)^tD_{v(X_r)}p_r(X_r)\bigr].
\]
For $a=b-D$, suppose $N_b\ge n^{-Q}$, where $Q$ is any exponent
fixed in advance.  Suppose also that either $p_b$ is a fixed
nondegenerate finite linear combination of concentrating visible
exact test gradients, or there is a subsequent block with
$N_{b+D}\le R N_b$ for a fixed finite $R$.
Then
\begin{equation}
 |A_a-e^{-\lambda D}A_b|
 +\bigl(\E\norm{p_a-\widetilde E A_a}^2\bigr)^{1/2}
 \le\eta e^{-\lambda D}N_b.
 \label{eq:cutoff-alignment}
\end{equation}
The ratio $R$ is used only to establish exactness at $b$; it does not
enter the choice of $D$ or the fluctuation estimate.
\end{proposition}

\begin{proof}
We may choose $D>2$.  First put the block at prescribed relative times
using \Cref{eq:cutoff-restart}.  In the terminal alternative take $\ell=D$
and $g=S_{T-b}f$.  In the subsequent-block alternative take $\ell=2D$
and $g=S_{T-b-D}f$.  In both cases the restart origin is $a=b-D$.
The new initial law is the old marginal at $a$, with the same upper
density bound and common trajectory Poincar\'e constant.  The gradient
and recipe contractions are unchanged at matching times; in particular
$\widehat A_r=A_{a+r}$.  The terminal alternative is the same
condition on $\widehat p_D$, while the other alternative gives
$\widehat N_{2D}\le R\widehat N_D$.
For the limiting argument denote the restarted process again by $X$
and use $a=0,b=D$; retain also $2D$ in the second alternative.
Normalize $N_D=1$.

We first establish exactness at $D$, then control the mean and
fluctuation of the slow coefficient, and finally recover a uniform
finite-dimensional estimate.  In the terminal alternative,
$q_b$ lies in the visible exact sector by assumption.  In the
subsequent-block alternative, the budget $N_{2D}\le R$ allows
\Cref{prop:exact-closedness} at a good time selected in a fixed
subinterval compactly contained in $(D,2D)$.  Both adjacent time
separations thus stay positive after extraction.  Transport back to the prescribed time $D$ by
\Cref{prop:exact-transport} gives the same conclusion at $b=D$.  The value of the finite budget $R$ has served
this purpose and is not used again.

Retain $0,D$ and the good time $s\in[D-1,D-1/2]$ in the same
scalar extraction, together with $2D$ and the later good time when
that alternative is used.  If $c_r$ denotes the random slow coefficient of
$q_r$, projection and \Cref{prop:exact-transport} give
\begin{equation}
 c_a=e^{-\lambda(s-a)}\E[c_s\mid X_a],\qquad
 \E c_a=e^{-\lambda(b-a)}\E c_b.
 \label{eq:cutoff-slow-transport}
\end{equation}
This transports the mean at the desired rate.  To align the vector
itself in unweighted square mean, we must also make the centered
coefficient small.

Choose a good time $s\in[b-1,b-1/2]$ as in
\Cref{eq:cutoff-good-time}.  Approximate one slow column by a coarse
smooth scaled recipe $B$.  The scalar function
\[
 F_s(x)=B(x)^tD_{v(x)}p_s(x)
\]
approximates the corresponding coefficient $c_s$ in scalar $L^2$
to error at most $C\delta$, where $\delta$ is the unweighted
approximation accuracy.  Indeed multiplication by $D_v$ is a contraction
in the ordinary norm, and \Cref{eq:exact-contraction-continuity} extends
the contraction with $p_s$ from fixed recipes to the slow column.
The bound uses typicality under the evolving law; stationary
mean-square approximation alone would not give it.

The trajectory Poincar\'e inequality will control this fluctuation once
we bound the gradient of $F_s$.  Set $c_i=v_iB_i$; the exact product rule
gives
\[
 d_jF_s=\sum_i c_i(x^j)d_jp_{s,i}
                   +\sum_i p_{s,i}(x)d_jc_i.
\]
The columns of $B$ have entries bounded by $C_B/\sqrt n$, and the
matrix $(d_jc_i)_{ij}$ has operator norm at most $C_B/\sqrt n$.
For the first term, mixed-gradient symmetry changes
$d_jp_{s,i}$ to $d_ip_{s,j}$.  Also
$v_i(x^j)\le e^{C|J_{ij}|}v_i(x)$ and $v_i\le2w_i$, so
\[
 \sum_i\frac{c_i(x^j)^2}{w_i(x)}\le C_B.
\]
Cauchy--Schwarz, followed by summation over $j$, proves
\begin{equation}
 \E\norm{dF_s}^2
 \le C_B\left(
   \E\sum_iw_i\norm{d_ip_s}^2+n^{-1}\E\norm{p_s}^2\right)
 \le C_B.
 \label{eq:cutoff-coefficient-derivative}
\end{equation}
The same estimate applies to every one of the finitely many columns.

Before choosing these recipes or $D$, freeze once and for all a
moderate-rate instance of \Cref{eq:cutoff-propagation}, say with
$\rho_0=\lambda/2$, and its constant $C_{\rho_0}$.  The Poincar\'e
bound at time $a$, this instance of propagation applied to the scalar
function $F_s$, and \Cref{eq:cutoff-coefficient-derivative} give
\begin{equation}
 \Var\bigl(\E[F_s(X_s)\mid X_a]\bigr)
 \le C_P C_{\rho_0}C_B e^{-2\rho_0(s-a)}+o(1).
 \label{eq:cutoff-coefficient-fluctuation}
\end{equation}
The coefficient approximation costs $C\delta$ in $L^2$ before
transport.  By \Cref{eq:cutoff-slow-transport}, both this error and
the square root of \Cref{eq:cutoff-coefficient-fluctuation} carry
the factor $e^{-\lambda(s-a)}$.  Relative to
$e^{-\lambda D}$, their cost is bounded by
\[
 C\delta+C(C_PC_{\rho_0}C_B)^{1/2}e^{-\rho_0(D-1)}.
\]
In particular the approximation error does not acquire a factor
$e^{\lambda D}$.

The hidden part and the exact fast complement, propagated from
$s$, cost at most
$C e^{-(\lambda+\kappa)(D-1)}$ for some $\kappa>0$, by
\Cref{thm:hidden-propagation,prop:exact-isolation,lem:exact-upgrade}.
The latter also applies when the visible vector is only in
$\mathcal H$ with weighted image in the fast part of $\mathcal X_v$.
After division by $e^{-\lambda D}$, this fast-sector error is bounded by
$C e^{\lambda+\kappa}e^{-\kappa D}$.
First choose $\delta$ sufficiently small in terms of $\eta$ and $C_E$,
and fix the coarse recipes and $C_B$.  Now choose $D$ large so that
both normalized exponential errors,
\[
 C(C_PC_{\rho_0}C_B)^{1/2}e^{-\rho_0(D-1)}
 \quad\text{and}\quad C e^{\lambda+\kappa}e^{-\kappa D},
\]
are smaller than the required fixed fraction of $\eta$.
Restoring the common factor $e^{-\lambda D}$ proves the limiting mean
and fluctuation estimate with exact slow coefficients and reconstruction.

The block length is now fixed, so we may choose the finer recipes
$\widetilde E$ without altering the fluctuation estimate.  At this
$D$, an accuracy smaller than
$c\eta e^{-(\lambda+C)D}$ makes all changes in the reported $A_a$,
$A_b$, and the reconstruction smaller than the remaining fraction
of $\eta e^{-\lambda D}$.  Here $e^{CD}$ is a rough bound for any
endpoint norm needed in this comparison.  The finer derivative
constants have not replaced $C_B$ in
\Cref{eq:cutoff-coefficient-fluctuation}.

Finally use strict slack, the deterministic good-disorder selection
in the proof of \Cref{prop:cutoff-block}, and normalization by $N_b$.
The condition $N_b\ge n^{-Q}$ preserves a fixed polynomial envelope;
choose every Poincar\'e and propagation remainder with an exponent
larger than $2Q+H+2$ before using it.  For a violating sequence, perform the one-block or two-block restart
above before extraction.  The normalized gradients, actual conditional
laws, common Poincar\'e bound, and endpoint ratio meet the same
hypotheses at the prescribed relative endpoints; any parameter selected from the family is
fixed conditional on the disorder before the path is sampled.
The resulting fixed-interval sequence would contradict the limiting
estimate.  Returning by \Cref{eq:cutoff-restart} proves
\Cref{eq:cutoff-alignment} uniformly at the original block positions,
without a rate for any fixed-time $o(1)$.  This includes the
parameter selected from the preparation family below.
\end{proof}

\subsection{Preparation, the cone, and a macroscopic mean}

Choose $m$ fixed scalar tests $F_1,\ldots,F_m$ from the potential
class defining $\mathcal X$ whose gradient projections onto the
slow space form an invertible matrix.  Such tests exist by density.
Replace them by fixed linear combinations so that the limiting
matrix is the identity.  Their gradients have uniformly bounded
pointwise Euclidean norm, and the functions have polynomial
supremum bounds.  For a unit output coefficient $u\in\R^m$, the
terminal gradient $d(\sum_j u_jF_j)$ consequently satisfies
\begin{equation}
 |A_T|\ge c_0,\qquad N_T\le K_0|A_T|
 \label{eq:cutoff-terminal-cone}
\end{equation}
uniformly on typically initialized trajectories, after the recipe
accuracy is sufficiently small.

The preparation must have a common Poincar\'e bound even when the slow
recipes are refined.  We establish this uniformity before making the
final approximation choices.  For any fixed smooth bounded scaled recipe matrix
$\widetilde E$ and any $0<\alpha<1/2$, set
\[
 R_n=n^{1/2-\alpha},\qquad
 g^z(x)=\widetilde E(x)z,\qquad |z|<R_n.
\]
Let $T_s^z$ be the rate-one refresh dynamics whose target mean at
the old configuration is
\[
 m_i^z(x)=\tanh((Jx)_i+g_i^z(x)).
\]
The target may depend on the spin being refreshed, so this process
need not be reversible.  Its generator is nevertheless exactly
\begin{equation}
 L^z f=\sum_i(m_i^z-x_i)d_if.
 \label{eq:cutoff-preparation-generator}
\end{equation}

\begin{lemma}[Common preparation bounds]
\label{lem:cutoff-preparation}
For each fixed recipe matrix and all sufficiently large $n$,
uniformly for $|z|<R_n$ and $0<h\le1$, the law
$\nu_z=\mu T_h^z$ satisfies
\begin{align}
 \log\norm{\dd\nu_z/\dd\mu}_\infty
   &\le nh(e^{2\varepsilon_n}-1)=O(n^{1-\alpha}),
   &\varepsilon_n&=\sup_{z,x,i}|g_i^z(x)|=O(n^{-\alpha}),
   \label{eq:cutoff-preparation-density}\\
 \Var_{\nu_z}f
   &\le\left[C_\mu e^{C_*h}
         +4\frac{e^{C_*h}-1}{C_*}\right]\nu_z\norm{df}^2.
   \label{eq:cutoff-preparation-poincare}
\end{align}
Here $C_\mu$ can be the inverse equilibrium gap, and $C_*$ is a
common constant depending on the original interaction bounds, not
on the fixed recipe.  A finer recipe changes the sufficiently large
dimension threshold, but not this common constant.
\end{lemma}

\begin{proof}
Both the added field and its difference matrix are small.  Scaled recipe
calculus gives
\begin{equation}
 \sup_{z,x}\norm{g^z(x)}_\infty+
 \sup_{z,x}\norm{(d_i g_j^z(x))_{ij}}_\op
 \le C_{\widetilde E}n^{-\alpha}.
 \label{eq:cutoff-small-perturbation}
\end{equation}
At each refresh attempt, including an attempt that leaves the spin
unchanged, the ratio of the perturbed outcome probability to the
ordinary outcome probability is at most $e^{2\varepsilon_n}$.
On the common attempt construction the path likelihood is therefore
at most $e^{2\varepsilon_n N_h}$, with
$N_h\sim\operatorname{Poisson}(nh)$ under the ordinary process.
Started at $\mu$, the ordinary endpoint has law $\mu$ conditional
on $N_h=k$, because its embedded update kernel preserves $\mu$.
Thus $X_h$ and $N_h$ are independent under this stationary reference
law, and for every endpoint $y$,
\[
 \frac{\mu T_h^z(y)}{\mu(y)}
 \le\E_\mu[e^{2\varepsilon_n N_h}\mid X_h=y]
 =\exp\{nh(e^{2\varepsilon_n}-1)\}.
\]
Averaging the attempt likelihood therefore proves the endpoint density
bound \Cref{eq:cutoff-preparation-density}.

For the variance estimate, we use gradients of the perturbed semigroup.
Differentiate
\Cref{eq:cutoff-preparation-generator}.  The exact gradient
extension has coefficients $a_{ij}^z=d_im_j^z$, including the
diagonal $a_{ii}^z$.  Its formula is
\[
 (L^z-I)p_i-2x_i\sum_j a_{ij}^z d_jp_i+\sum_j a_{ij}^z p_j.
\]
The diagonal $a_{ii}^z p_i$ is an $o(1)$ zero-order term; the
own-spin derivative of an exact gradient is zero.  Taylor expansion
of $\tanh$, with \Cref{eq:cutoff-small-perturbation} and the
bounded row and column square sums, gives common bounds
\[
 \norm{a^z}_\op\le C,\qquad
 \max_i\sum_j\frac{|a_{ij}^z|^2}{w_j^z}\le C,
 \qquad w_j^z=1-x_jm_j^z.
\]
For example the leading matrix is the field-difference matrix times
the diagonal variance; $v_j^z\le2w_j^z$ controls the displayed
weighted row sum.  The quadratic Taylor remainders are bounded
Schur products of matrices with bounded row and column square sums.
Once the right side of \Cref{eq:cutoff-small-perturbation} is at most
one, the square-evolution proof of \Cref{prop:path-rough} therefore
has a common constant $C_*$.  In particular
\begin{equation}
 \norm{dT_s^zf}^2\le e^{C_*s}T_s^z\norm{df}^2.
 \label{eq:cutoff-preparation-gradient}
\end{equation}
Thus the gradient estimate remains valid for this possibly nonreversible
preparation: it was derived from the stated generator and its exact
gradients, without requiring a Gibbs invariant law.

Its jump bracket is at most $4\norm{dT_{h-s}^zf}^2$.
The conditional martingale variance formula and
\Cref{eq:cutoff-preparation-gradient} give
\[
 \Var_x(f(X_h))
 \le4\frac{e^{C_*h}-1}{C_*}\,T_h^z\norm{df}^2(x).
\]
The other term in total variance is
\[
 \Var_\mu(T_h^zf)
 \le C_\mu\mu\norm{dT_h^zf}^2
 \le C_\mu e^{C_*h}\mu T_h^z\norm{df}^2.
\]
Their sum is \Cref{eq:cutoff-preparation-poincare}.
\end{proof}

We now fix the remaining parameters in an order that preserves the
common constants.  First freeze the
moderate rate $\rho_0=\lambda/2$ and its propagation constants in
\Cref{eq:cutoff-propagation}.  Use
\Cref{lem:cutoff-preparation} with $h\le1$ in
\Cref{eq:cutoff-trajectory-poincare}; this gives a common $C_P$
for every prepared ordinary trajectory, independent of later recipe
refinements.  Next fix the output tests, choose
$K>\max\{K_0,2C_E+2\}$, and choose $\eta$ small in terms of $K$
and the desired rate slack.  Choose the coarse coefficient recipes,
then the alignment block $D$, and then the finer reporting and
preparation recipes, in exactly the order in the proof of
\Cref{prop:cutoff-coefficients}.  Finally choose a positive $h\le1$
for these fine recipes.  Shrinking $h$ cannot increase the common
Poincar\'e constant; larger fine-recipe derivative constants change
the dimension threshold in \Cref{lem:cutoff-preparation} and the
choice of $h$, not the previously fixed coefficient fluctuation
bound or $D$.

\begin{lemma}[Iteration inside the slow cone]
\label{lem:cutoff-cone}
Given $\xi_1>0$ and $\epsilon_0>0$, the preceding choices can be
made so that, for $t$ an integer multiple of $D$, the gradient
$p_0=dS_t(\sum_j u_jF_j)$ under any prepared law satisfies
\begin{equation}
 |A_0|\ge c e^{-(\lambda+\xi_1)t},\qquad
 \bigl(\E_{\nu_z}\norm{p_0-\widetilde E A_0}^2\bigr)^{1/2}
       \le\epsilon_0|A_0|,
 \label{eq:cutoff-cone-output}
\end{equation}
uniformly for $|u|=1$, $|z|<R_n$, and positive times
$D\le t\le A\log n$.
\end{lemma}

\begin{proof}
The cone compares the full gradient norm with its mean slow coefficient.
Assume $N_b\le K|A_b|$ at the end of a block.
\Cref{eq:cutoff-alignment} gives
\begin{align*}
 |A_a|&\ge(1-K\eta)e^{-\lambda D}|A_b|,\\
 N_a&\le C_E'|A_a|+\eta e^{-\lambda D}N_b
       \le\left(C_E'+\frac{K\eta}{1-K\eta}\right)|A_a|,
\end{align*}
where $C_E'$ is arbitrarily close to $C_E$ after fine approximation.
Choose $K\eta<1/4$ and then $\eta$ small enough that the last
coefficient is at most $K$, and
$K\eta/(1-K\eta)\le\epsilon_0$.  Thus the cone is preserved.
At the terminal time it holds by \Cref{eq:cutoff-terminal-cone}.
At each step, apply the fixed-interval restart in the proof of
\Cref{prop:cutoff-coefficients} to the ordinary trajectory with
initial block law $\nu_zS_a$.  Its upper density bound follows from
\Cref{eq:cutoff-preparation-density,eq:cutoff-restart-density}, and
its common Poincar\'e bound is the one fixed above.  This restarts
only the ordinary part of the evolution, leaving the preparation
$T_h^z$ unchanged.  At subsequent blocks the required forward norm
ratio follows from the preceding cone step:
\[
 \frac{N_{b+D}}{N_b}
 \le\frac{K C_Ae^{\lambda D}}{1-K\eta},
 \qquad |A_b|\le C_A N_b.
\]
This ratio is finite for the already chosen block $D$.  Its only role
in \Cref{prop:cutoff-coefficients} is to establish exactness at the block
endpoint, so its dependence on $D$ does not create a circular choice.

After $k=t/D$ steps,
\[
 |A_0|\ge c_0(1-K\eta)^k e^{-\lambda t}
 =c_0e^{-[\lambda-D^{-1}\log(1-K\eta)]t}.
\]
Take $D\ge1$ and $\eta\le\xi_1/(2K)$, so that the extra exponent
is at most $\xi_1$.  To verify the inverse-power cutoff throughout
the iteration, choose
\[
 Q>A(\lambda+\xi_1)+2.
\]
The lower bound for $|A|$, and $|A|\le C_A N$, then keep $N$ above
$n^{-Q}$ for all the $O(\log n)$ blocks.  Choose the inverse-power
accuracies used in the supporting Poincar\'e estimates larger than
$2Q+H+2$.  No new union bound over fixed-time limiting errors is
needed: \Cref{eq:cutoff-alignment} holds on one event uniformly over
the functions, parameter values, and comparison times in question.
The residual estimate in \Cref{eq:cutoff-cone-output} follows from
the final cone step.
\end{proof}

\begin{lemma}[Preparation creates a ball of final means]
\label{lem:cutoff-jacobian}
For sufficiently small fixed $h>0$ chosen after the fine recipes,
the smooth map
\[
 \mathcal M_t(z)=\bigl(\mu T_h^z S_tF_j\bigr)_{j=1}^m,
 \qquad |z|<R_n,
\]
has Jacobian with smallest singular value at least
\begin{equation}
 \sigma_{\min}(D\mathcal M_t(z))
       \ge c h e^{-(\lambda+\xi_1)t}=:a_t.
 \label{eq:cutoff-jacobian}
\end{equation}
Its image therefore contains the open Euclidean ball of radius
$a_tR_n$ centered at $(\mu F_j)_{j=1}^m$.
\end{lemma}

\begin{proof}
We estimate the Jacobian in each unit output direction.  For $|u|=1$,
put $F=\sum_j u_jF_j$ and apply \Cref{eq:cutoff-cone-output} at the exit
from preparation.  Differentiation on the finite state space and
Duhamel's formula give the row vector
\begin{equation}
 u^tD\mathcal M_t(z)
 =\int_0^h\E_{\mu T_s^z}
   \left[(dT_{h-s}^zS_tF)^tD_{v^z}\widetilde E\right]\dd s,
 \qquad v_i^z=1-(m_i^z)^2.
 \label{eq:cutoff-duhamel-jacobian}
\end{equation}
Indeed $\partial_{z_k}L^zf=
\sum_i v_i^z\widetilde E_{ik}d_if$, also when the target depends on
the old spin.

Write $p_0=\widetilde E A_0+r$ under $\nu_z$.  The rough vector
extension behind \Cref{eq:cutoff-preparation-gradient} propagates
$r$ back through preparation at cost at most $e^{C_*h/2}$ in
mean square.  Its contribution to
\Cref{eq:cutoff-duhamel-jacobian} is therefore
$O(h\epsilon_0|A_0|)$.

For the main vector, uniformly over $0\le s\le h$,
\begin{equation}
 \left(\E_{\mu T_s^z}
   \norm{K_{h-s}^z\widetilde E-\widetilde E}^2\right)^{1/2}
       \le\omega_{\widetilde E}(h)+o_n(1),
 \qquad \omega_{\widetilde E}(h)\longrightarrow0.
 \label{eq:cutoff-short-continuity}
\end{equation}
Here the norm is the sum over the fixed number of columns.  To
verify this estimate, the diagonal likelihood representation has
row moments bounded by $e^{C h}$, and its likelihood difference
from one has squared mean $O(h)$, by the likelihood martingale
bracket and the bounded row square sums.  The entries of a fixed
scaled recipe are bounded by $C/\sqrt n$, so this diagonal error
costs $O(h)$ after summation.  The part with at least one inserted
coefficient matrix costs $O(h)$ in norm by its Duhamel expansion
and the common rough bound.  Finally the ordinary recipe value
changes by $o_h(1)$ in spatial mean square along the short path:
its smooth Euclidean extension has Lipschitz constant
$C_{\widetilde E}/\sqrt n$, while the expected number of changed
spins is at most $nh$.  These estimates also apply to the perturbed
rates, using \Cref{eq:cutoff-small-perturbation}; the own-spin
zero-order correction is $o_n(1)$.  They prove
\Cref{eq:cutoff-short-continuity} without a regularity assertion for
the exact eigenvectors.

By \Cref{eq:cutoff-preparation-density}, every marginal in this
calculation has upper relative log-density $o(n)$.
The one-state typicality comparisons, the chosen fine accuracy,
and $\norm{v^z-v}_\infty=o(1)$ therefore give
\[
 \E_{\mu T_s^z}[\widetilde E^tD_{v^z}\widetilde E]
       =I+O(\delta_{\rm fine})+o_n(1)
\]
uniformly in $s,z$.  Combining this with
\Cref{eq:cutoff-duhamel-jacobian,eq:cutoff-short-continuity} yields
\[
 \left|u^tD\mathcal M_t(z)-hA_0^t\right|
 \le Ch\bigl(\epsilon_0+\delta_{\rm fine}
                   +\omega_{\widetilde E}(h)+o_n(1)\bigr)|A_0|.
\]
Choose $\epsilon_0$ and the fine accuracy first, and then choose a
positive fixed $h$ so that the factor in parentheses is small.
\Cref{eq:cutoff-cone-output} now proves
\Cref{eq:cutoff-jacobian}, uniformly for every unit output vector.

To pass from the Jacobian estimate to a ball of attainable means, we
lift straight line segments by the inverse Jacobian.  This argument
requires only the proved singular-value bound, not global injectivity.
For a target vector $v$ with $|v|<a_tR_n$, solve the inverse-Jacobian
ordinary differential equation
\[
 z'(r)=(D\mathcal M_t(z(r)))^{-1}v,\qquad z(0)=0,
 \qquad 0\le r\le1.
\]
The right side is smooth and $|z'(r)|\le |v|/a_t<R_n$.
The solution stays in a compact subball of the parameter ball and
therefore exists to time one.  Differentiation gives
$\mathcal M_t(z(r))=\mathcal M_t(0)+rv$.
At $z=0$ the preparation is ordinary stationary evolution, so
$\mathcal M_t(0)=(\mu F_j)_j$.  This proves the assertion.
\end{proof}

\begin{lemma}[Lower cutoff for an isolated visible rate]
\label{lem:cutoff-visible-lower}
If $\lambda<\lambda_{\rm sp}$, the lower conclusion of
\Cref{lem:cutoff-spin-lower} still holds.
\end{lemma}

\begin{proof}
Fix $0<\epsilon<1$, take $0<\alpha<\epsilon/8$ and
$0<\xi_1<\lambda\epsilon/4$, and round
$(1-\epsilon)t_n$ \emph{up} to an integer multiple $t$ of $D$.
The radius in \Cref{lem:cutoff-jacobian} is at least
\[
 c h e^{-(\lambda+\xi_1)D}
 n^{\,1/2-\alpha-(1-\epsilon)(\lambda+\xi_1)/(2\lambda)}.
\]
The exponent is positive, so the ball of attainable means has diverging
radius.  In particular, one may choose a preparation for which the mean
of $F_1$ differs from $\mu F_1$ by a quantity tending to infinity.
The variances of $F_1$ under $\nu_zS_t$ and $\mu$ are bounded:
use the common trajectory Poincar\'e estimate and its bounded
pointwise gradient for the first, and the equilibrium gap for the
second.  Chebyshev's inequality implies
$\norm{\nu_zS_t-\mu}_{\TV}\to1$.
Since $\nu_z$ is a probability law on starting configurations,
convexity bounds this distance by
$\sup_x\norm{S_t(x,\cdot)-\mu}_{\TV}$.
Monotonicity then gives the same lower bound at the unrounded,
earlier time $(1-\epsilon)t_n$.
\end{proof}

\begin{theorem}[Continuous-time cutoff]
\label{thm:cutoff-continuous}
For every fixed $0<\beta<1$ and $\epsilon\in(0,1)$, put
$t_n=\log n/(2\lambda)$.  Then
\begin{equation}
 \begin{aligned}
 \sup_x\norm{S_{(1-\epsilon)t_n}(x,\cdot)-\mu}_{\TV}
       &\longrightarrow1,\\
 \sup_x\norm{S_{(1+\epsilon)t_n}(x,\cdot)-\mu}_{\TV}
       &\longrightarrow0,
 \end{aligned}
 \label{eq:cutoff-continuous}
\end{equation}
in probability over the disorder.  The number $\lambda>0$ is given
entirely by the dimension-first equilibrium construction in
\Cref{eq:exact-lambda}.
\end{theorem}

\begin{proof}
The upper bound follows from \Cref{lem:cutoff-upper}.  For the lower
bound, \Cref{lem:cutoff-spin-lower} applies when
$\lambda=\lambda_{\rm sp}$, and \Cref{lem:cutoff-visible-lower} applies
when $\lambda<\lambda_{\rm sp}$.  These cases exhaust the definition
of $\lambda$.
\end{proof}

\Cref{prop:model-independent} supplies the continuous-time conclusion
at $\beta=0$, with $\lambda(0)=1$.  It remains to transfer both cases to
the discrete update clock.

\subsection{The discrete update clock}

\begin{proposition}[Cutoff for single-site update attempts]
\label{prop:cutoff-clock}
Let $P=I+L/n$, so one step selects a uniform site and refreshes its
spin, including attempts that leave the state unchanged.  For every
fixed $0\le\beta<1$, its worst-case total-variation cutoff is at
\[
 \frac{n\log n}{2\lambda(\beta)},
\]
in probability over the disorder.
\end{proposition}

\begin{proof}
We prove both directions of clock transfer.  Write $d_c(t)$ and
$d_d(k)$ for the continuous and discrete distances; their kernels satisfy
the exact Poissonization identity $S_t=e^{nt(P-I)}$.  For
$k=\lfloor(1-\epsilon)nt_n\rfloor$ and
$T=(1-\epsilon/2)t_n$, convexity and discrete-time contraction give
\[
 d_c(T)\le\P\{\operatorname{Poisson}(nT)<k\}+d_d(k).
\]
The probability is $O((nt_n)^{-1})$, while
$d_c(T)\to1$, proving the discrete lower bound.

{\emergencystretch=3em
For the upper bound, a uniform holding probability permits the local
Poisson-to-binomial lemma.  Restrict to the high-probability event
$\norm J_\op\le K$.  At every configuration,
$n^{-1}\sum_i(Jx)_i^2\le K^2$, so at least half of the fields have
absolute value at most $\sqrt2K$.  The probability of keeping a
selected spin is at least $(1+e^{2|(Jx)_i|})^{-1}$.  Consequently\par}
\begin{equation}
 P(x,x)\ge h_0:=\frac1{2(1+e^{2\sqrt2K})}>0.
 \label{eq:cutoff-holding}
\end{equation}
Choose $\vartheta=1-h_0/2\in(0,1)$ and write
\[
 P=(1-\vartheta)I+\vartheta R,
 \qquad S_t=e^{\vartheta nt(R-I)}.
\]
Here $R$ is a Markov kernel preserving $\mu$.
Set $M=\lceil(1+\epsilon)nt_n\rceil$.  For every $x$ and $|f|\le1$,
put $u_j=R^jf(x)-\mu f$, so $|u_j|\le2$.  For each fixed real $y$,
\[
 \E u_{\operatorname{Poisson}(\vartheta M+y\sqrt M)}
 =S_{M/n+y\sqrt M/(\vartheta n)}f(x)-\mu f.
\]
The time correction is $O_y(\sqrt{\log n/n})$, and hence this
time exceeds $(1+\epsilon/2)t_n$ for large $n$.  The absolute value
of the display is bounded by
$2d_c((1+\epsilon/2)t_n)\to0$, uniformly over $x,f$.

We have verified every shifted Poisson average required by
\Cref{lem:input-clock}, uniformly over starts and tests.  Apply that
uniform transfer: its assumptions are a fixed
$\vartheta\in(0,1)$, $M\to\infty$, bounded arrays, and vanishing
Poisson averages for every fixed shift $y$, all just verified.
It gives
\[
 \E u_{\operatorname{Binomial}(M,\vartheta)}
       =P^Mf(x)-\mu f\longrightarrow0
\]
uniformly over the test family.  Taking half its supremum proves
$d_d(M)\to0$.  For convergence in disorder probability, the proof of
that lemma uses a fixed finite signed approximation by shifted
Gaussian densities; its finitely many Poisson averages are all
bounded by the same random quantity $2d_c((1+\epsilon/2)t_n)$.
Thus it applies on the present high-probability events without any
additional convergence rate.
\end{proof}
